% 01-intro.tex — §1 Introduction
\section{Introduction}

The execution layer of a blockchain information system is a replicated
consensus state machine: every node on the consensus path must produce
bit-identical results for the same transaction. The EVM is a widely
deployed smart-contract execution environment; when a client enables a
JIT, the compiler becomes part of the consensus-critical path.
Peephole rewrites in the JIT act directly during machine-code generation,
so their correctness directly affects the determinism of on-chain
execution; an incorrect equivalence judgment, deployed on only some
clients, may trigger a chain fork.

Existing formal verification of EVM peephole rewrites mostly targets
the block-level bytecode abstraction. The machine-level patterns
emitted by a JIT during code generation---especially the multi-word
carry chains that arise when the 256-bit word size is lowered onto
64-bit hardware---sit at a different abstraction layer from bytecode
rules and are awkward to validate as standard LLVM/Alive2 peephole
queries. The root cause is how IRs model bit-level side results. In
LLVM IR, basic arithmetic is defined only on $N$-bit
integers; the carry bit must be retrieved through an intrinsic that
returns a multi-return tuple of \{result, carry\}. Rules depending on
this carry bit can be expressed only through indirect structures in
Alive~\cite{alive} and similar tools; the bottleneck is the IR's
expressiveness, not the SMT solver~\cite{demoura2008z3}. For a
consensus-critical JIT, an unchecked rule class leaves the engineer
with three choices---drop it (performance), ship it without verification
(safety), or hand-prove it (scalability).

Our approach is to extend the IR with first-class bit-vector operators so
that SMT-checkable peephole rules cover a larger rule class, without
modifying the solver or relaxing the differential-test threshold. In
dMIR, the bit-level patterns that previously required indirect
multi-return encoding are represented directly, so the
corresponding rules fall within Z3 bit-vector equivalence. The SMT scope
is dMIR arithmetic/flag semantics; gas, memory, storage, and x86
equivalence remain outside it.

Concretely, we make three contributions:
\begin{itemize}
\item \textbf{Carry-aware dMIR semantics.} We promote carry and overflow
from implicit side results to first-class dMIR operators, so
ADC/SBB-style multi-word dependencies become single bit-vector
equivalence obligations rather than indirect multi-return encodings.
\item \textbf{SMT admission and production coverage.} An offline Z3
pipeline admits all \textbf{70} production dMIR rewrites under
fixed-width modular bit-vector semantics. Its enumeration also reaches
all production left-hand sides (\textbf{70/70}), including carry-chain
forms that resist direct formulation in LLVM IR.
\item \textbf{Layered optimization with consensus-path evidence.}
SMT-checked dMIR arithmetic simplifications compose with x86-64 CgIR
cleanups validated by matched differential testing; across 27
evmone-bench workloads, the 83-rule system attains a \textbf{+7.24\%}
geometric-mean speedup and passes \textbf{5{,}884} Cancun differential
tests with byte-identical final state and matching gas accounting within
suite coverage.
\end{itemize}
